\begin{thebibliography}{99}
\raggedright
\bibitem{ABDP} M. A. Alekseyev, W. Bruns, J. Dong and S. Palcoux,
\emph{Classifying integral Grothendieck rings up to rank 5 and beyond},
\href{https://arxiv.org/abs/2507.07023v2}{arXiv:2507.07023v2} (2025).

\bibitem{BNRW} P. Bruillard, S.-H. Ng, E. C. Rowell and Z. Wang,
\emph{Rank-finiteness for modular categories}, J. Amer. Math. Soc.
\textbf{29} (2016), 857--881;
\href{https://arxiv.org/abs/1310.7050v3}{arXiv:1310.7050v3}.
\bibitem{DP} J. Dong and S. Palcoux,
\emph{Rank-four Grothendieck rings through multiplicity 32},
manuscript and exact-certificate supplement, 17 September 2026.
The numbered results cited here refer to that manuscript.
\bibitem{HR} S.-M. Hong and E. C. Rowell,
\emph{On the classification of the Grothendieck rings of non-self-dual
modular categories}, J. Algebra \textbf{324} (2010), 1000--1015;
\href{https://arxiv.org/abs/0907.1051}{arXiv:0907.1051}.
\bibitem{Izumi} M. Izumi,
\emph{A Cuntz algebra approach to the classification of near-group categories},
in \emph{Proceedings of the 2014 Maui and 2015 Qinhuangdao Conferences
in Honour of Vaughan F. R. Jones' 60th Birthday},
Proc. Centre Math. Appl. Austral. Nat. Univ. \textbf{46} (2017), 222--343;
\href{https://arxiv.org/abs/1512.04288v2}{arXiv:1512.04288v2}.
\bibitem{NS} S.-H. Ng and P. Schauenburg,
\emph{Frobenius--Schur indicators and exponents of spherical categories},
Adv. Math. \textbf{211} (2007), 34--71;
\href{https://arxiv.org/abs/math/0601012}{arXiv:math/0601012}.
\bibitem{LPR} Z. Liu, S. Palcoux and Y. Ren,
\emph{Classification of Grothendieck rings of complex fusion categories
of multiplicity one up to rank six}, Lett. Math. Phys. \textbf{112}
(2022), article 54;
\href{https://arxiv.org/abs/2010.10264v5}{arXiv:2010.10264v5}.
\bibitem{OEIS} OEIS Foundation Inc.,
\emph{The On-Line Encyclopedia of Integer Sequences},
\href{https://oeis.org/A354473}{A354473}, accessed 18 September 2026.
\bibitem{VS} G. Vercleyen and J. K. Slingerland,
\emph{On low rank fusion rings}, J. Math. Phys. \textbf{64} (2023),
article 091703;
\href{https://arxiv.org/abs/2205.15637v4}{arXiv:2205.15637v4}.
Input: the ancillary file \texttt{FusionRingMultiplicationTables}.

\bibitem{DGNO} V. Drinfeld, S. Gelaki, D. Nikshych and V. Ostrik,
\emph{On braided fusion categories I}, Selecta Math. (N.S.)
\textbf{16} (2010), 1--119;
\href{https://arxiv.org/abs/0906.0620v3}{arXiv:0906.0620v3}.
\bibitem{Ostrik} V. Ostrik,
\emph{Pivotal fusion categories of rank 3}, Mosc. Math. J.
\textbf{15} (2015), no.~2, 373--396;
\href{https://arxiv.org/abs/1309.4822v1}{arXiv:1309.4822v1}.
\bibitem{Thornton} J. Thornton, \emph{On braided near-group categories},
\href{https://arxiv.org/abs/1102.4640v1}{arXiv:1102.4640v1} (2011).
\end{thebibliography}
